\documentclass[a4paper,12pt]{article}
\usepackage[utf8]{inputenc}
\usepackage[polish]{babel}
\usepackage[T1]{fontenc}
\usepackage{geometry}
\usepackage{graphicx}
\usepackage{enumitem}
\usepackage{hyperref}
\usepackage{longtable}
\geometry{margin=2.5cm}

\title{Specyfikacja Projektu Zaliczenowego (Programowanie w Języku Java)}
\author{Bartosz Badura, Eryk Witkowski}
\date{[Data]}

\begin{document}

\maketitle

\section*{1. Tytuł Projektu}

\vspace{0.5 cm}
 \noindent
\textbf{Analizator ruchu sieciowego}

\section*{2. Opis i Cele Projektu}
\textit{Określ, jaki problem rozwiązuje projekt, do czego służy oraz jakie są jego główne cele funkcjonalne.}

\vspace{0.5 cm}
 \noindent Projekt ma na celu stworzenie aplikacji, która umożliwi monitorowanie ruchu sieciowego oraz wykrywanie podejrzanych aktywności. Główne cele projektu:

\begin{itemize}
    \item Wyświetlanie ruchu w sieci w czasie rzeczywistym.
    \item Możliwość podglądu pojedynczych pakietów.
    \item Wykrywanie oraz ogłaszanie podejrzanych zachowań.
    \item Generowanie raportów o anomaliach.
    \item Pamiętanie zaobserwowanych nieprawidłowości.
\end{itemize}

\section*{3. Wymagania Funkcjonalne}
\textit{Lista konkretnych funkcji, które aplikacja powinna oferować użytkownikowi końcowemu.}

\vspace{0.5 cm}
 \noindent
\underline{Przykład:} 

\begin{itemize}
    \item Przegląd ruchu.
    \item Podgląd pojedynczych pakietów.
    \item Zapisywanie wybranych pakietów.
    \item Wykrywanie spamu.
    \item Wykrywanie ataku DDOS.
    \item Wykrywanie niepoprawnych przesyłów.
    \item Generowanie raportów.
\end{itemize}

\section*{4. Wymagania Niefunkcjonalne}
\textit{Wymagania techniczne, jakościowe i pozafunkcjonalne, które aplikacja powinna spełniać.}

\vspace{0.5 cm}
 \noindent
\underline{Przykład:} 
\begin{itemize}
    \item \textbf{Wydajność:} Czas przetwarzania danych poniżej ... sekundy.
    \item \textbf{Skalowalność:} Obsługa minimum ... użytkowników równocześnie.
    \item \textbf{Bezpieczeństwo:} HTTPS, hasła szyfrowane (bcrypt).
    \item \textbf{Użyteczność:} Interfejs zgodny z WCAG (AA).
    \item \textbf{Zgodność:} Responsywny interfejs (desktop, tablet, mobilne).
\end{itemize}

\section*{5. Architektura Systemu}
\textit{Opis architektury systemu, jego głównych komponentów oraz zależności między nimi.}

\vspace{0.5 cm}
 \noindent
\underline{Przykład:} 
Typ architektury: \textbf{Klient-Serwer (MVC)}.

\subsection*{Schemat:}
\begin{itemize}
    \item Frontend: React.js + TailwindCSS
    \item Backend: Java (framework zależny od implementacji)
    \item Baza danych: PostgreSQL lub MongoDB
    \item WebSocket: np. STOMP/WebSocket
    \item Autoryzacja: JWT / OAuth2
\end{itemize}


\section*{6. Zastosowane Algorytmy}
\textit{Opis algorytmów i struktur danych wykorzystanych w aplikacji do realizacji funkcjonalności.}

\vspace{0.5 cm}
 \noindent
\underline{Przykład:} 
\begin{itemize}
    \item \textbf{Algorytm sortowania zadań:} Sortowanie zadań po terminie, priorytecie lub statusie — wykorzystanie np. sortowania przez porównanie z wykorzystaniem Comparator w Javie.

    \item \textbf{Algorytmy uwierzytelniania:} Tokenizacja i walidacja JWT (algorytm HMAC-SHA256), szyfrowanie haseł z użyciem bcrypt.

    \item \textbf{Filtrowanie i wyszukiwanie:} Przeszukiwanie listy zadań według etykiet, tekstu lub użytkownika — implementacja z wykorzystaniem filtrów i indeksowania.

    \item \textbf{Zarządzanie sesją:} Wygasanie tokenów na podstawie czasu (timestamp + TTL), middleware do odświeżania sesji.

    \item \textbf{Algorytmy notyfikacji:} Sprawdzanie zbliżających się terminów i wyzwalanie powiadomień.

    \item \textbf{Algorytmy kolejkowania wiadomości:} W WebSocket można zastosować buforowanie i kolejkowanie aktualizacji do wielu klientów jednocześnie.
\end{itemize}

\section*{7. Struktura Klas}
\textit{Opisuje strukturę klas, ich odpowiedzialności oraz podział logiczny na warstwy w aplikacji.}

\vspace{0.5 cm}
 \noindent
\underline{Przykład:} 
\vspace{0.5 cm}
 \noindent
\subsection*{Model}
\begin{itemize}
    \item \texttt{User} – identyfikator, imię, email, hasło, rola.
    \item \texttt{Task} – id, tytuł, opis, termin, status, przypisany użytkownik.
    \item \texttt{Tag} – id, nazwa, kolor, powiązanie z zadaniami.
\end{itemize}

\subsection*{Serwisy}
\begin{itemize}
    \item \texttt{UserService} – rejestracja, logowanie, zarządzanie rolami.
    \item \texttt{TaskService} – logika zarządzania zadaniami.
    \item \texttt{NotificationService} – generowanie powiadomień.
    \item \texttt{AuthService} – JWT, walidacja, bezpieczeństwo.
\end{itemize}

\subsection*{Kontrolery}
\begin{itemize}
    \item \texttt{UserController} – REST API użytkownika.
    \item \texttt{TaskController} – REST API zadań.
    \item \texttt{AuthController} – logowanie i uwierzytelnianie.
\end{itemize}

\subsection*{Dodatkowe komponenty}
\begin{itemize}
    \item \texttt{WebSocketHandler} – aktualizacje w czasie rzeczywistym.
    \item \texttt{DTO} – obiekty wymiany danych (Data Transfer Object).
    \item \texttt{Repository} – interfejsy do bazy danych (np. JPA/Hibernate).
\end{itemize}

\section*{8. Wykorzystane Technologie i Biblioteki}
\textit{Opis technologii oraz narzędzi użytych w projekcie.}

\vspace{0.5 cm}
 \noindent
\underline{Przykład:} 
\vspace{0.5 cm}
 \noindent
\subsection*{Frontend:}
React, React Query

\subsection*{Backend:}
Java, WebSocket

\subsection*{Baza Danych:}
PostgreSQL lub MongoDB

\subsection*{DevOps:}
Docker, GitHub Actions

\section*{9. API i Modele Danych}
\textit{Przykładowe punkty końcowe API oraz struktura danych wymienianych między systemami.}

\vspace{0.5 cm}
 \noindent
\underline{Przykład:} 
\subsection*{Przykładowe Endpointy:}
\begin{verbatim}
GET /api/tasks
POST /api/tasks
PUT /api/tasks/:id
DELETE /api/tasks/:id
\end{verbatim}

\subsection*{Model Danych:}
\begin{verbatim}
{
  "id": "uuid",
  "tytul": "string",
  "opis": "string",
  "termin": "timestamp",
  "priorytet": "enum(niski, sredni, wysoki)",
  "etykiety": ["string"],
  "zrealizowane": "boolean",
  "przypisanyDo": "user_id"
}
\end{verbatim}

\section*{10. Interfejs Użytkownika}
\textit{Opis interfejsu graficznego aplikacji, jego struktury, logiki działania oraz dostępnych ekranów.}

\vspace{0.5 cm}
 \noindent
\underline{Przykład:} 
\subsection*{Główne widoki:}
\begin{itemize}
    \item \textbf{Ekran logowania i rejestracji:} Formularze z walidacją danych, możliwość odzyskania hasła.
    \item \textbf{Pulpit użytkownika:} Przegląd aktualnych zadań, kalendarz, szybki dostęp do przypisanych zadań.
    \item \textbf{Zarządzanie zadaniami:} Widok listy zadań z możliwością sortowania, filtrowania, edytowania i usuwania.
    \item \textbf{Tworzenie i edycja zadań:} Formularz z polami na tytuł, opis, etykiety, termin, przypisanie użytkownika.
    \item \textbf{Powiadomienia:} Wyskakujące alerty o nadchodzących terminach i zmianach w zadaniach (real-time).
    \item \textbf{Panel administracyjny:} Zarządzanie użytkownikami, rolami, dostępem.
\end{itemize}

\subsection*{Elementy UX/UI:}
\begin{itemize}
    \item Responsywność — dostosowanie do urządzeń mobilnych.
    \item Konsekwentna kolorystyka, nawigacja boczna lub górna.
    \item Łatwość obsługi — maksymalnie 3 kliknięcia do najważniejszych funkcji.
    \item Tryb ciemny/jasny — dostępny z poziomu ustawień.
\end{itemize}

\subsection*{Technologie UI:}
React, TailwindCSS, React Router, Headless UI, ikony Lucide.

\section*{11. Etapy Realizacji / Harmonogram}
\textit{Plan wdrożenia projektu z podziałem na fazy i szacowanym czasem realizacji.}

\vspace{0.5 cm}
 \noindent
\underline{Przykład:} 
\begin{longtable}{|c|p{8cm}|c|}
\hline
\textbf{Faza} & \textbf{Opis} & \textbf{Czas} \\
\hline
Planowanie & Określenie wymagań, wybór technologii & 1 tydzień \\
\hline
Backend & Budowa API, modele danych & 2 tygodnie \\
\hline
Frontend & Interfejs, integracja z API & 2 tygodnie \\
\hline
Testy i realtime & WebSocket, testy jednostkowe i integracyjne & 1 tydzień \\
\hline
Wdrożenie & Konfiguracja Docker, CI/CD, serwer & 1 tydzień \\
\hline
\end{longtable}

\section*{12. Testowanie}
\textit{Sposoby testowania aplikacji, typy testów oraz używane narzędzia.}

\vspace{0.5 cm}
 \noindent
\underline{Przykład:} 
\begin{itemize}
    \item Testy jednostkowe (logika backendu i frontend).
    \item Testy integracyjne (API).
    \item Testy E2E – np. Cypress, Selenium.
    \item Testy obciążeniowe – np. JMeter, k6.
\end{itemize}

\section*{13. Ryzyka i Środki Zaradcze}
\textit{Identyfikacja potencjalnych zagrożeń projektowych i strategii ich ograniczenia.}

\vspace{0.5 cm}
 \noindent
\underline{Przykład:} 
\begin{longtable}{|p{5cm}|p{3cm}|p{6cm}|}
\hline
\textbf{Ryzyko} & \textbf{Wpływ} & \textbf{Zarządzanie} \\
\hline
Nadmierna liczba funkcji & Wysoki & Zablokowanie wymagań po planowaniu \\
\hline
Błędy bezpieczeństwa & Wysoki & Audyty kodu, sprawdzone biblioteki \\
\hline
Opóźnienia & Średni & Praca iteracyjna (Agile) \\
\hline
\end{longtable}

\section*{14. Kryteria Sukcesu}
\textit{Opisz po czym poznać, że projekt zakończył się sukcesem.}

\vspace{0.5 cm}
 \noindent
\underline{Przykład:} 
\vspace{0.5 cm}
 \noindent
\begin{itemize}
    \item Wszystkie funkcje działają zgodnie ze specyfikacją.
    \item Pokrycie testami > 80\%.
    \item Brak błędów krytycznych na produkcji.
    \item Pozytywna opinia użytkowników (średnia ocena \textgreater 4.5).
\end{itemize}

\section*{15. Szacunkowe Nakłady Pracy (Workload Estimation)}
\textit{Osobne oszacowanie czasu pracy dla poszczególnych modułów i funkcji.}

\vspace{0.5 cm}
 \noindent
\underline{Przykład:} 
\vspace{0.5 cm}
 \noindent
Całkowity szacowany czas pracy: \textbf{160 godzin roboczych}.

\begin{longtable}{|p{6cm}|p{4cm}|p{3cm}|}
\hline
\textbf{Moduł / Zadanie} & \textbf{Opis} & \textbf{Szac. czas (h)} \\
\hline
Analiza i projektowanie & Wymagania, mockupy, wybór technologii & 12 \\
\hline
Backend API (CRUD, JWT) & Modele danych, logika biznesowa & 35 \\
\hline
Frontend UI (React) & Komponenty, interakcje z API & 30 \\
\hline
Uwierzytelnianie i sesje & Login, rejestracja, autoryzacja & 15 \\
\hline
WebSocket – realtime & Notyfikacje i live updates & 10 \\
\hline
Testy (unit + E2E) & Pokrycie testami kluczowych ścieżek & 15 \\
\hline
Docker / CI / Deployment & Konfiguracja środowiska, produkcja & 10 \\
\hline
Dokumentacja techniczna & README, API schema, komentarze & 8 \\
\hline
Bufor czasowy / błędy & Debugowanie, nieprzewidziane problemy & 25 \\
\hline
\end{longtable}

\end{document}

